\section{Реализация интеллектуальной системы}
\label{sec:development}

\subsection{Выбор средств реализации}

\textbf{Языки программирования, фреймворки и библиотеки}

% Требования:
% - перечислить используемые языки, фреймворки, библиотеки;
% - указать версии основных средств.

\textbf{Обоснование выбора средств}

% Требования:
% - обосновать выбор с учётом типа интерфейса, требований к надёжности, производительности и интеграции;
% - указать, как выбранные средства поддерживают расширяемость.

\textbf{Реализация пользовательского интерфейса}

% Требования:
% - описать соответствие реализованных экранов спроектированным макетам;
% - указать механизмы обработки ввода, навигации, отображения результатов и ошибок.

\subsection{Реализация программного (или аппаратного) интерфейса}

% Требования:
% - описать реализацию основных операций API/протоколов;
% - привести примеры запросов/ответов, логов обмена или временных диаграмм.

\subsection{Тестирование интерфейса}

\textbf{Функциональное и интеграционное тестирование}

% Требования:
% - описать перечень тестовых сценариев, покрывающих ключевые функции;
% - указать результаты тестирования, найденные ошибки и способы их устранения.

\textbf{Оценка удобства и качества взаимодействия}

% Требования:
% - оценить соответствие интерфейса требованиям usability и «интеллектуальности»;
% - при наличии указать методику оценки (опрос, экспертная оценка, эвристики).

\subsection{Вывод по разделу реализации}

% Требования:
% - перечислить реализованные сценарии и поддерживаемые типы взаимодействия;
% - указать ограничения текущего прототипа и возможные направления доработки;
% - подчеркнуть готовность интерфейса к демонстрации в режиме реальной работы.
